\documentclass[10pt,letterpaper]{article}

% ---------- packages ----------
\usepackage[utf8]{inputenc}
\usepackage[T1]{fontenc}
\usepackage{microtype}
\usepackage{amsmath,amssymb}
\usepackage{graphicx}
\usepackage{booktabs}
\usepackage{multirow}
\usepackage{xcolor}
\usepackage{natbib}
\usepackage{hyperref}
\usepackage[capitalise]{cleveref}
\usepackage{enumitem}
\usepackage[margin=1in]{geometry}
\usepackage{xspace}

\hypersetup{
    colorlinks=true,
    linkcolor=blue!60!black,
    citecolor=blue!60!black,
    urlcolor=blue!60!black
}

\newcommand{\todo}[1]{\textcolor{red}{\textbf{[TODO: #1]}}}
\newcommand{\method}{\textsc{hippoVLA}\xspace}
\newcommand{\sTMB}{\textsc{sTMB}\xspace}

% ---------- title ----------
\title{\method: Hippocampus-Inspired Spatiotemporal\\
Disentanglement for Vision-Language-Action Models}

\author{
    Anonymous Authors\\
    \texttt{\{anonymous\}@institution}
}

\date{}

\begin{document}
\maketitle

% =====================================================================
\begin{abstract}
Vision-Language-Action (VLA) models inherit the monolithic transformer backbone of their VLM ancestors, asking a single context window to simultaneously serve two cognitively distinct functions: per-frame \emph{spatial grounding} and across-frame \emph{temporal integration}. We argue that this entanglement is a structural cause of two persistent VLA failures --- weak long-horizon execution and brittle spatial reasoning --- because the two functions compete for the same finite representational capacity.
We draw on the Complementary Learning Systems (CLS) account of human memory, in which the neocortex slowly consolidates statistical regularities while the hippocampus rapidly encodes episodic context, and propose \method, a brain-inspired VLA that explicitly disentangles the two streams. \method couples a spatially enhanced VLM backbone (analogous to dorsal-stream cortex) with a Short-Term Memory Bank (\sTMB), a hippocampus-inspired episodic memory featuring (i) pattern-separated learnable slots, (ii) sparse \emph{theta-rhythm} temporal sampling, and (iii) a novel \emph{spatial-saliency-aware gate} that selectively writes only spatially relevant tokens into memory. On LIBERO 4-in-1 \method reaches $99.8/94.4/96.4/96.0$ on Object/Spatial/Goal/Long, and on the CALVIN ABC$\rightarrow$D long-horizon leaderboard it attains an average rollout length of $4.38$ (rank 3), a $+0.68$ absolute gain over its no-memory ablation. Ablations confirm that neither the spatial backbone nor the memory module suffices alone: the two streams are \emph{synergistic}, mirroring the CLS division of labour.
\end{abstract}

% =====================================================================
\section{Introduction}
\label{sec:intro}

Vision-Language-Action (VLA) models extend Vision-Language Models (VLMs) to embodied control by adding an action head on top of a multimodal transformer~\citep{black2024pi0,kim2024openvla,brohan2023rt2}. While this recipe has produced impressive single-step manipulation, it has consistently struggled at two ends of the difficulty spectrum: \emph{long-horizon} tasks that require integrating information across many frames, and \emph{spatially demanding} tasks that require precise 3D grounding. We identify a structural cause that has so far been under-examined: the same self-attention stack is asked to do both, and the two functions trade off against one another inside a fixed context window.

\paragraph{Spatial vs.\ temporal in one stream.}
A frame-level VLA token serves two purposes simultaneously: it grounds language onto the current scene, and it is the only carrier of past visual history into the future. As context grows, attention dilutes spatial detail; as the model is tuned for spatial precision, long-range temporal patterns become harder to retain. We empirically observe both failure modes on LIBERO and CALVIN: vanilla strong-spatial backbones cap out below $90\%$ on long-horizon tasks (\Cref{tab:libero-cross}), and methods optimised for long horizons concede spatial accuracy.

\paragraph{A brain-inspired alternative.}
Biological agents do not solve this tension with a single homogeneous network. The Complementary Learning Systems (CLS) hypothesis~\citep{mcclelland1995cls,kumaran2016cls} posits two interacting memories: a \emph{slow} neocortical system that learns generalisable representations across many experiences, and a \emph{fast} hippocampal system that rapidly encodes individual episodes and replays them during consolidation and retrieval. The hippocampus is not a passive buffer: dentate-gyrus pattern separation produces sparse, decorrelated memory traces; theta-gamma coupling imposes a discrete temporal rhythm on encoding; and prefrontal feedback gates which incoming events are worth committing to memory at all~\citep{rolls2013mechanisms}.

\paragraph{This work.}
We propose \method, a VLA architecture that mirrors this division of labour. A spatially enhanced backbone (RynnBrain-8B) plays the role of dorsal-stream/neocortex, providing rich per-frame spatial grounding. A new module --- the \textbf{Short-Term Memory Bank} (\sTMB) --- plays the role of hippocampus, maintaining a fixed-size pool of pattern-separated slots that absorb information across frames at a sparse temporal rhythm. Crucially, \sTMB writes are governed by a \emph{spatial-saliency-aware gate} that consults the spatial backbone to decide which incoming tokens are worth memorising. This makes the two streams \emph{coupled} rather than merely concatenated: the spatial system informs what the temporal system stores.

\paragraph{Contributions.}
\begin{itemize}[leftmargin=*]
    \item We diagnose a structural spatial--temporal entanglement in monolithic VLA backbones and propose a CLS-inspired disentangled design.
    \item We introduce \sTMB, a hippocampus-inspired memory module with pattern-separated slots, sparse temporal sampling, and a novel spatial-saliency-aware write gate.
    \item On LIBERO 4-in-1 \method achieves $99.8/94.4/96.4/96.0$ across the four suites, and on CALVIN ABC$\rightarrow$D it attains an average rollout length of $4.38$ (rank 3), a $+0.68$ absolute gain over the same backbone without memory.
    \item Ablations show that the spatial backbone and memory module are \emph{individually insufficient but jointly sufficient}, supporting the CLS interpretation: the two streams are synergistic, not redundant.
\end{itemize}

% =====================================================================
\section{Related Work}
\label{sec:related}

\paragraph{Vision-Language-Action models.}
Recent VLA systems --- OpenVLA~\citep{kim2024openvla}, $\pi_0$~\citep{black2024pi0}, RT-2~\citep{brohan2023rt2}, GR00T~\citep{nvidia2025gr00t} --- attach an action decoder (discrete tokens, MLP regression, or diffusion) to a pretrained VLM. Most treat the VLM as a black-box feature extractor and do not explicitly model temporal context beyond the native context window.

\paragraph{Memory in robot learning.}
Recurrent and transformer-based policies use implicit memory through hidden states or context~\citep{brohan2023rt2,reed2022gato}. Explicit memory has been explored as external retrieval~\citep{wu2022memorizing}, slot attention~\citep{locatello2020slot}, and memory tokens~\citep{bulatov2022rmt}, but rarely with an explicit gate that consults a separate spatial stream. Concurrent VLA work such as \todo{cite recent memory-augmented VLAs} adds context windows but does not factorise spatial vs.\ temporal capacity.

\paragraph{Brain-inspired AI.}
Hippocampal models in deep learning include differentiable neural computers~\citep{graves2016dnc}, neural episodic control~\citep{pritzel2017nec}, and CLS-style replay~\citep{kumaran2016cls,parisi2019cls}. Our work differs in two ways: it instantiates CLS \emph{architecturally} (two parallel streams), not just procedurally (replay buffer), and the gate is conditioned on a sister network rather than self-conditioned.

\paragraph{Spatial foundation models.}
Spatially enhanced VLMs --- e.g.\ SpatialBot, RynnBrain --- improve geometric reasoning by training on point/depth-grounded data. We adopt RynnBrain-8B as the dorsal-stream component but, crucially, we do not merely fine-tune it: we couple it to a memory module that consumes its spatial-saliency signal.

% =====================================================================
\section{Method}
\label{sec:method}

\method consists of three components (\Cref{fig:arch}): a spatial stream (\Cref{sec:spatial}), the hippocampus-inspired \sTMB (\Cref{sec:stmb}), and the action decoder. We describe each, and then state how they integrate.

\begin{figure}[t]
\centering
\todo{Figure: two-stream architecture diagram. Left: spatial stream (RynnBrain-8B) producing per-frame visual tokens with a spatial saliency map. Right: \sTMB with $N$ slots, theta-sampled write path, spatial-saliency-aware gate. Bottom: action decoder reading from both streams.}
\caption{\method architecture. The spatial stream (left) produces per-frame tokens and a saliency signal. \sTMB (right) maintains a fixed pool of pattern-separated slots updated at a sparse temporal rhythm; its write gate is conditioned on the spatial saliency, so only spatially relevant tokens are committed to memory. The action decoder reads from both streams, mirroring the CLS interaction between hippocampus and neocortex.}
\label{fig:arch}
\end{figure}

\subsection{Cognitive framing}
\label{sec:framing}
We treat a VLA forward pass at time $t$ as a function $a_t = \pi(o_t, h_t)$ where $o_t$ is the current observation and $h_t$ summarises history. In a monolithic transformer, $h_t$ is implicit in the past tokens still inside the context window, and the same parameters that compute $h_t$ also compute the spatial features of $o_t$. Following CLS, we factor $\pi$ into a fast-episodic stream $\mathcal{H}$ (hippocampus) and a slow-semantic stream $\mathcal{N}$ (neocortex):
\begin{equation}
    a_t = D\bigl(\,\mathcal{N}(o_t),\; \mathcal{H}(o_{\le t})\,\bigr),
\end{equation}
where $D$ is a lightweight action decoder. The two streams are \emph{coupled}: $\mathcal{H}$ writes to memory only with the consent of $\mathcal{N}$, via a spatial-saliency-aware gate (\Cref{sec:gate}).

\subsection{Spatial Stream: RynnBrain-8B}
\label{sec:spatial}
We use RynnBrain-8B as the spatial backbone. For each frame $o_t$ it produces visual tokens $V_t \in \mathbb{R}^{N_v \times d}$ aligned with language tokens. RynnBrain is pretrained with point/depth-grounded supervision so that $V_t$ carries an implicit 3D prior; we additionally extract a per-token spatial saliency $s_t \in [0,1]^{N_v}$ defined as the L2 norm of the residual between RynnBrain's spatially grounded path and a non-spatial baseline path \todo{finalise definition; alternative: attention rollout to the [PT] token of \textsc{ChainOfPoint}}. The saliency $s_t$ is the only signal that crosses from $\mathcal{N}$ to $\mathcal{H}$.

\subsection{\sTMB: Hippocampus-Inspired Episodic Memory}
\label{sec:stmb}

\paragraph{Memory bank as pattern-separated slots.}
\sTMB maintains a fixed-size memory $M \in \mathbb{R}^{B \times K \times d}$ with $K=64$ learnable slot embeddings, initialised i.i.d.\ Gaussian and persisted across frames within an episode. The slot count is small relative to per-frame token counts ($K \ll N_v$), enforcing a bottleneck analogous to dentate-gyrus pattern separation: the model must compress episodic content into distinct slots rather than store raw frames.

\paragraph{Theta-rhythm sparse sampling.}
Hippocampal encoding is gated by theta--gamma coupling: only a fraction of incoming sensory frames are committed to long-term memory. We mirror this with a tunable \emph{sampling interval} $\tau$: \sTMB writes once every $\tau$ frames rather than every frame. Empirically (\Cref{sec:exp-ablation}) $\tau=10$ outperforms $\tau=5$, consistent with the view that sparser, more decorrelated samples make better episodic anchors.

\paragraph{Write path with spatial-saliency-aware gate.}
\label{sec:gate}
At a write step $t$, slots attend to incoming visual tokens via cross-attention,
\begin{equation}
    \tilde M_t = \mathrm{CrossAttn}(M_{t-1},\, V_t + e_t),
\end{equation}
where $e_t$ is a learned temporal position embedding. A standard gated update would compute $g = \sigma(W_g[M_{t-1};\tilde M_t])$ and blend $M_t = g \odot \tilde M_t + (1{-}g)\odot M_{t-1}$. We extend this with a saliency-conditioned multiplicative term:
\begin{equation}
    g_t \;=\; \sigma\!\bigl(W_g\, [M_{t-1};\,\tilde M_t]\bigr) \;\odot\; \phi(\bar s_t),\qquad
    M_t = g_t \odot \tilde M_t + (1 - g_t) \odot M_{t-1},
    \label{eq:gate}
\end{equation}
where $\bar s_t \in [0,1]$ is the mean spatial saliency of $V_t$ and $\phi$ is a small MLP returning a per-slot scalar. The intuition is direct: when the spatial stream reports that the current frame is spatially uninformative (camera transit, occlusion, no manipulable object), the gate clamps shut and \sTMB is preserved; when the frame is spatially salient (object newly visible, gripper near contact), the gate opens. This is the architectural analog of prefrontal--hippocampal gating in CLS.

\paragraph{Read path.}
For action prediction at time $t$, the action decoder query tokens $Q_t$ retrieve from \sTMB via cross-attention,
\begin{equation}
    \hat Q_t = Q_t + \mathrm{CrossAttn}(Q_t, M_t).
\end{equation}
The decoder thus consumes a fused representation that contains \emph{both} the current spatial grounding (from $\mathcal{N}$) and the episodic context (from $\mathcal{H}$).

\subsection{Training and inference}
\sTMB is trained jointly with the action head; the spatial backbone is fine-tuned with a small learning rate. We initialise $M$ at episode start from the learnable prototype and reset between episodes. \sTMB adds $\sim$\todo{X}M parameters over the backbone, $<$\todo{X}\% of the total.

% =====================================================================
\section{Experiments}
\label{sec:exp}

We evaluate \method on two complementary embodied benchmarks: \textbf{LIBERO 4-in-1}~\citep{liu2023libero} (single co-trained policy across Spatial/Object/Goal/Long suites) and \textbf{CALVIN ABC$\rightarrow$D}~\citep{mees2022calvin} (long-horizon language-conditioned manipulation, primary metric: average completed instructions out of 5).

\subsection{Setup}
\paragraph{Training.} Co-training on LIBERO 4-in-1 with batch size 16, 20k--60k optimisation steps, AdamW. Backbone: starVLA-Qwen3VL-4B for the memory ablations, RynnBrain-8B-OFT for the cross-backbone comparison and CALVIN. \sTMB uses $K=64$ slots, $\tau \in \{5,10\}$, memory length 5.
\paragraph{Inference.} \sTMB is reset at episode start; memory persists across the full episode.

\subsection{LIBERO 4-in-1: memory ablation}
\label{sec:exp-libero-ablation}
\Cref{tab:libero-mem} reports success rates under the same Qwen3VL-4B backbone with different memory configurations. Adding \sTMB with $\tau=10$ improves Long by $+3.0$ and Spatial by $+1.2$ over the no-memory baseline, and matches the 50k-step official checkpoint with only 20k steps. The $\tau=5$ vs.\ $\tau=10$ comparison is consistent with the theta-rhythm interpretation: sparser, more decorrelated writes produce better episodic memory.

\begin{table}[t]
\centering
\small
\caption{LIBERO 4-in-1 memory ablation (success rate \%). Same Qwen3VL-4B backbone; only the memory configuration changes. \sTMB with $\tau=10$ is best on three of four suites.}
\label{tab:libero-mem}
\begin{tabular}{lcccccccc}
\toprule
Setting & BS & Steps & Backbone & Module & Object & Spatial & Goal & Long \\
\midrule
Official        & --- & 50k & Qwen3VL-4B & ---       & 99.6 & ---  & 99.0 & 98.6 \\
Reproduced      & --- & 50k & Qwen3VL-4B & ---       & 99.6 & 93.2 & 98.6 & 93.0 \\
+ \sTMB ($\tau{=}5$)  & 16  & 20k & Qwen3VL-4B & len 5 & 99.6 & 94.2 & 98.2 & 95.8 \\
+ \sTMB ($\tau{=}10$) \textbf{[ours]} & 16 & 20k & Qwen3VL-4B & len 5 & \textbf{99.8} & \textbf{94.4} & 96.4 & \textbf{96.0} \\
\bottomrule
\end{tabular}
\end{table}

\subsection{LIBERO: cross-backbone comparison}
\label{sec:exp-libero-cross}
\Cref{tab:libero-cross} shows that adding \sTMB on top of the RynnBrain-8B-OFT backbone improves Spatial by $+1.0$ while preserving the strong Object/Goal/Long performance of the spatial backbone. This is the central empirical signature of the CLS hypothesis: the spatial stream alone is bottlenecked on Spatial despite being a strong backbone, and the memory stream alone (on a weaker backbone) is bottlenecked on absolute ceiling --- only the union saturates both axes.

\begin{table}[t]
\centering
\small
\caption{LIBERO cross-backbone comparison. \sTMB lifts the spatially strong RynnBrain backbone above its no-memory ceiling on Spatial.}
\label{tab:libero-cross}
\begin{tabular}{lccccc}
\toprule
Method & Note & Spatial & Object & Goal & Long \\
\midrule
Qwen2.5-VL-GR00T-LIBERO-4in1            & local repro      & 98.4 & 99.0 & 96.2 & 91.0 \\
RynnBrain-8B-OFT (no memory)            & step 60k         & 94.4 & \textbf{100.0} & \textbf{98.0} & \textbf{95.8} \\
\textbf{\method} (RynnBrain + \sTMB)    & step 60k         & \textbf{95.4} & 99.6 & 97.2 & 93.8 \\
\bottomrule
\end{tabular}
\end{table}

\subsection{CALVIN ABC$\rightarrow$D long-horizon}
\label{sec:exp-calvin}
\Cref{tab:calvin} reports the full CALVIN long-horizon leaderboard. \method (RynnBrain + \sTMB) attains average rollout length $4.38$, ranking 3rd among published methods --- and crucially, $+0.68$ absolute over the no-memory ablation of the same backbone (3.70). This is the largest single contribution of any component in our pipeline. Long-horizon execution is precisely the regime in which monolithic transformers run out of context; \sTMB sidesteps this by maintaining episodic state outside the language context entirely.

\begin{table}[t]
\centering
\small
\caption{CALVIN ABC$\rightarrow$D long-horizon leaderboard. Avg.\ Len.\ is the mean number of consecutive instructions completed (max 5). \method ranks 3rd and is $+0.68$ over its own no-memory ablation.}
\label{tab:calvin}
\begin{tabular}{rlcccccc c}
\toprule
\# & Method & Step 1 & Step 2 & Step 3 & Step 4 & Step 5 & Avg.\ Len. \\
\midrule
1 & FLOWER                          & 99.4 & 95.8 & 90.7 & 84.9 & 77.8 & 4.53 \\
2 & UniVLA                          & 98.9 & 94.8 & 89.0 & 82.8 & 75.1 & 4.41 \\
3 & \textbf{\method (ours)}         & 99.1 & 94.4 & 88.4 & 81.6 & 74.1 & \textbf{4.38} \\
4 & SeeR-Large                      & 96.3 & 91.6 & 86.1 & 80.3 & 74.0 & 4.28 \\
5 & GR-MG                           & 96.8 & 89.3 & 81.5 & 72.7 & 64.4 & 4.04 \\
6 & MoDE                            & 96.2 & 88.9 & 81.1 & 71.8 & 63.5 & 4.01 \\
7 & RynnBrain (no memory, ours)     & 88.7 & 79.9 & 72.7 & 67.3 & 61.0 & 3.70 \\
8 & GHIL-Glue                       & 95.2 & 88.5 & 73.2 & 62.5 & 49.8 & 3.69 \\
9 & RoboUniView                     & 94.2 & 84.2 & 73.4 & 62.2 & 50.7 & 3.64 \\
10 & Diffusion Transformer Policy   & 94.5 & 82.5 & 72.8 & 61.3 & 50.0 & 3.61 \\
11 & CLOVER                         & 96.0 & 83.5 & 70.8 & 57.5 & 45.4 & 3.53 \\
12 & 3D Diffuser Actor              & 92.2 & 78.7 & 63.9 & 51.2 & 41.2 & 3.27 \\
13 & GR-1                           & 85.4 & 71.2 & 59.6 & 49.7 & 40.1 & 3.06 \\
14 & DeeR                           & 86.2 & 70.1 & 51.8 & 41.5 & 30.4 & 2.82 \\
15 & SuSIE                          & 87.0 & 69.0 & 49.0 & 38.0 & 26.0 & 2.69 \\
16 & RoboFlamingo                   & 82.4 & 61.9 & 46.6 & 33.1 & 23.5 & 2.47 \\
17 & SPIL                           & 74.2 & 46.3 & 27.6 & 14.7 &  8.0 & 1.71 \\
18 & HULC                           & 41.8 & 16.5 &  5.7 &  1.9 &  1.1 & 0.67 \\
\bottomrule
\end{tabular}
\end{table}

\subsection{Ablations}
\label{sec:exp-ablation}

\paragraph{(a) Streams in isolation vs.\ combined.} We compare four configurations: (i) backbone only, (ii) backbone + \sTMB without saliency gate (i.e.\ vanilla GRU-style gate), (iii) backbone + \sTMB with saliency gate, and (iv) saliency gate but \emph{shuffled} across frames as a control. \todo{Run and fill table; expected: (iii) > (ii) > (i), and (iv) collapses to (i), confirming the saliency signal carries content.}

\paragraph{(b) Sampling interval $\tau$.} \Cref{tab:libero-mem} shows $\tau=10$ beats $\tau=5$. \todo{Add $\tau \in \{1, 3, 5, 10, 20\}$ scan to plot a clean curve.}

\paragraph{(c) Number of slots $K$.} \todo{$K \in \{8, 16, 32, 64, 128\}$ ablation; expected dome shape with peak around $K=64$ (pattern-separation argument predicts a sweet spot).}

\paragraph{(d) Slot probing.} \todo{Probe individual memory slots after rollout; show that some slots specialise to object identity and others to gripper--object spatial relations, paralleling place vs.\ object cells.}

\subsection{Qualitative analysis}
\Cref{fig:rollouts} shows representative rollouts on the four LIBERO suites. \todo{Add figure with 4 frame strips. Highlight cases where the no-memory baseline loses object identity after occlusion but \method recovers it from \sTMB.}

\begin{figure}[t]
\centering
\todo{Figure: 4 rollout strips (Spatial / Object / Goal / Long) with frame thumbnails and an indicator of when \sTMB writes are gated open vs.\ closed.}
\caption{Qualitative rollouts of \method on LIBERO. The bottom bar indicates the spatial-saliency gate state: open (filled) when the current frame is committed to \sTMB, closed (hatched) otherwise.}
\label{fig:rollouts}
\end{figure}

% =====================================================================
\section{Discussion}
\label{sec:discussion}

\paragraph{Why does the saliency gate matter?}
A simple data-driven gate (sigmoid over $[M;\tilde M]$) can in principle learn the same behaviour, but it must do so from sparse action-prediction gradient. Conditioning on $s_t$ injects an inductive bias: ``trust the spatial stream's judgment of what is worth remembering.'' This is precisely the role attributed to prefrontal--hippocampal feedback in mammalian memory.

\paragraph{Limitations.}
\sTMB is short-term: it does not consolidate across episodes. CLS predicts a second, slower process that distils episodic content into the cortex; we leave this to future work (e.g.\ replay-based fine-tuning of the backbone). The saliency definition is heuristic; a learned saliency head trained jointly with action prediction is a natural extension.

\paragraph{Reusing pretrained spatial backbones.}
A natural concern is that \method merely repackages an existing spatial backbone (RynnBrain). \Cref{tab:libero-cross} and the no-memory ablation in \Cref{tab:calvin} show this is not the case: the backbone alone, even fine-tuned, is bounded; the gain comes from the \emph{coupling} between the spatial stream and \sTMB through the saliency gate.

% =====================================================================
\section{Conclusion}

We presented \method, a hippocampus-inspired Vision-Language-Action model that disentangles spatial grounding from temporal integration via a Complementary Learning Systems-style two-stream architecture. The proposed Short-Term Memory Bank, with its pattern-separated slots, sparse theta-rhythm sampling, and spatial-saliency-aware gate, lifts a strong spatial backbone above its no-memory ceiling on both LIBERO and CALVIN. We see this as a step toward VLA systems whose architecture --- not merely their training data --- reflects the cognitive structure of the embodied agents they are meant to emulate.

% =====================================================================
\small
\bibliographystyle{plainnat}
\begin{thebibliography}{99}

\bibitem[Black et al.(2024)]{black2024pi0}
K.~Black et al.\ \newblock $\pi_0$: A vision-language-action flow model for general robot control. 2024.

\bibitem[Brohan et al.(2023)]{brohan2023rt2}
A.~Brohan et al.\ \newblock RT-2: Vision-language-action models transfer web knowledge to robotic control. 2023.

\bibitem[Bulatov et al.(2022)]{bulatov2022rmt}
A.~Bulatov et al.\ \newblock Recurrent memory transformer. \emph{NeurIPS}, 2022.

\bibitem[Graves et al.(2016)]{graves2016dnc}
A.~Graves et al.\ \newblock Hybrid computing using a neural network with dynamic external memory. \emph{Nature}, 2016.

\bibitem[Kim et al.(2024)]{kim2024openvla}
M.~Kim et al.\ \newblock OpenVLA: An open-source vision-language-action model. 2024.

\bibitem[Kumaran et al.(2016)]{kumaran2016cls}
D.~Kumaran, D.~Hassabis, J.~L.~McClelland. \newblock What learning systems do intelligent agents need? Complementary learning systems theory updated. \emph{Trends in Cognitive Sciences}, 2016.

\bibitem[Liu et al.(2023)]{liu2023libero}
B.~Liu et al.\ \newblock LIBERO: Benchmarking knowledge transfer for lifelong robot learning. 2023.

\bibitem[Locatello et al.(2020)]{locatello2020slot}
F.~Locatello et al.\ \newblock Object-centric learning with slot attention. \emph{NeurIPS}, 2020.

\bibitem[McClelland et al.(1995)]{mcclelland1995cls}
J.~L.~McClelland, B.~L.~McNaughton, R.~C.~O'Reilly. \newblock Why there are complementary learning systems in the hippocampus and neocortex. \emph{Psychological Review}, 1995.

\bibitem[Mees et al.(2022)]{mees2022calvin}
O.~Mees et al.\ \newblock CALVIN: A benchmark for language-conditioned policy learning for long-horizon robot manipulation tasks. \emph{IEEE RA-L}, 2022.

\bibitem[NVIDIA(2025)]{nvidia2025gr00t}
NVIDIA. \newblock GR00T N1: An open foundation model for generalist humanoid robots. 2025.

\bibitem[Parisi et al.(2019)]{parisi2019cls}
G.~I.~Parisi et al.\ \newblock Continual lifelong learning with neural networks: A review. \emph{Neural Networks}, 2019.

\bibitem[Pritzel et al.(2017)]{pritzel2017nec}
A.~Pritzel et al.\ \newblock Neural episodic control. \emph{ICML}, 2017.

\bibitem[Reed et al.(2022)]{reed2022gato}
S.~Reed et al.\ \newblock A generalist agent. \emph{TMLR}, 2022.

\bibitem[Rolls(2013)]{rolls2013mechanisms}
E.~T.~Rolls. \newblock The mechanisms for pattern completion and pattern separation in the hippocampus. \emph{Frontiers in Systems Neuroscience}, 2013.

\bibitem[Wu et al.(2022)]{wu2022memorizing}
Y.~Wu et al.\ \newblock Memorizing transformers. \emph{ICLR}, 2022.

\end{thebibliography}

\end{document}
